\chapter{Conclusions and Future Work}
\label{chapter:conclusions_and_future_work}


\section{Conclusions}




\section{Future Work}

\subsection{Integration of Meteorological Radar Data for Enhanced Rainfall Spatialization}

While the current data pipeline successfully spatializes point-based rain gauge measurements using elevation-based 
drift (e.g., Kriging with External Drift), this methodology presents inherent limitations when modeling extreme weather 
events. Elevation acts as a static covariate that effectively represents long-term, climatological precipitation trends 
(the orographic effect). However, it cannot capture the highly dynamic, localized, and heterogeneous nature of real-time 
convective storm cells. Consequently, the interpolated rainfall tensor may misrepresent the exact geographical 
distribution of the precipitation nucleus during a specific simulation step.

To achieve professional-grade meteorological realism and improve the predictive accuracy of the hydraulic simulation, 
a primary line of future work is the integration of Meteorological Radar Data (such as the C-band and S-band radar 
networks operated by the Spanish Meteorological Agency, AEMET).

This integration would fundamentally upgrade the Rainfall Engine through the following technical implementations:

\begin{itemize}
	\item \textbf{Radar Reflectivity Conversion}: Radar systems provide high-resolution spatial matrices of atmospheric 
	reflectivity ($Z$). The pipeline would need to implement algorithms to dynamically fetch these grids and convert 
	them into estimated rainfall intensity rates ($R$) using empirical formulas, such as the Marshall-Palmer relationship 
	($Z = aR^b$).
	\item \textbf{Radar-Rain Gauge Merging}: Because radar data is susceptible to systemic anomalies (e.g., beam blockage, 
	signal attenuation, or overestimation due to hail), it cannot be used in isolation. The future pipeline should 
	implement advanced merging techniques—such as Conditional Merging (CM) or Regression Kriging. In this upgraded 
	architecture, the radar imagery would serve as the highly detailed spatial covariate, while the physical rain gauges 
	would provide the "ground-truth" anchor points to correct the radar bias.
	\item \textbf{Enhanced Tensor Fidelity}: By replacing a static topographic covariate with a dynamic meteorological 
	one, the resulting $Time \times Y \times X$ data cube would precisely reflect the trajectory, shape, and intensity 
	of storm events.
\end{itemize}

Implementing this feature would drastically reduce spatial interpolation errors and provide the Cellular Automata 
($m:n-CA^k$) core with a highly realistic input tensor, which is strictly necessary for the accurate modeling of 
rapid-onset flash floods.
